\section{Empirical Evaluation}

\subsection{Performance Metrics}
We evaluate the algorithms based on the following metrics:

\begin{itemize}
    \item \textbf{Caption quality}: At the testing phase, we evaluate the generated difference captions using standard metrics, including BLEU-4 \citep{papineni2002bleu}, METEOR \citep{banerjee2005meteor}, ROUGE-L \citep{lin2004rouge}, and CIDEr-D \citep{vedantam2015cider}. These metrics measure how well the generated descriptions align with human annotations.
\end{itemize}

\subsection{Baseline Methods}
We consider the following recent state-of-the-art method as our primary baseline:

\begin{itemize}
    \item \textbf{CLIP4IDC} \citep{guo2022clip4idc}: A framework that leverages the pre-trained vision-language model, CLIP, to transfer cross-modal knowledge for image difference captioning. It aligns fine-grained visual differences with textual semantics to effectively capture and describe changes between image pairs.
\end{itemize}

CLIP4IDC is a strong baseline in IDC and represents a competitive representation learning framework for bridging visual differences and natural language descriptions.

\subsection{Benchmark Tasks or Datasets}
We evaluate the algorithms on several benchmark datasets widely used in the image difference captioning (IDC) literature:

\begin{itemize}
    \item \textbf{CLEVR-Change} \citep{park2019robust}: A synthetic dataset for scene-level difference captioning, where the model must identify semantic changes (e.g., object color, shape, and position) while ignoring non-semantic variations such as viewpoint and illumination changes.

    \item \textbf{Spot-the-Diff} \citep{jhamtani2018learning}: A real-world image-pair dataset that requires the model to generate natural language descriptions of semantic differences between two images.

    \item \textbf{Image-Editing-Request} \citep{tan2019expressing}: A dataset consisting of image pairs and corresponding editing instructions, where the model needs to describe fine-grained and diverse editing operations between images.
\end{itemize}
